\section{Lecture 2 (September 3rd)}
\begin{thm}
(Born rule for density operators) Using the spectral decomposition of a certain state ${A}=\sum a|a\rangle \langle a|=\sum a\mathrm{Pr}_{a}$, we can express the probability of observing ${A}$ and obtaining $a$ as 
\[\mathop{\mathrm{Tr}}({\rho }\mathop{\mathrm{Pr}_{a}})\]
\end{thm}
\vspace{2ex}
\begin{thm}
(Von Neumann equation) We check that the density matrix satisfies, for a pure state,
\[\begin{align*}
\dfrac{d \rho }{d t}=&\sum _{i}\Big(p_{i}|\psi _{i}\rangle \langle \psi _{i}|\Big)\\
=&\sum _{i}p_{i}\dfrac{d }{d t}|\psi _{i}\rangle \langle \psi _{i} |+\sum _{i}p_{i}|\psi _{i}\rangle \dfrac{d }{d t}\langle \psi_{i}|\\
=&\sum _{i}p_{i}\Big(\dfrac{H}{i\hbar }|\psi _{i}\rangle \Big)\langle \psi _{i}|+\sum _{i}p_{i}|\psi _{i}\rangle \Big(\dfrac{H}{-i\hbar }\langle \psi _{i}\Big)\\
=&\dfrac{H}{i\hbar }\sum _{i}p_{i}|\psi _{i}\rangle \langle \psi _{i}|-\sum p_{i}|\psi _{i}\rangle \langle \psi _{i}|\dfrac{H}{i\hbar }\\
=&\dfrac{1}{i\hbar }[H,\rho ]=-\dfrac{i}{\hbar }[H,\rho ]
\end{align*}\]
We see that mixed states and pure states evolve under the same Hamiltonian law.
\end{thm}
\vspace{2ex}
\begin{rmk}
(Necessary condition for a pure state) Notice that the square of the density matrix is equal to
\[{\rho }^2=\sum _{j}p_{j}^2|\psi _{j}\rangle \langle \psi _{j}|\]
Therefore, in general, we say mixed states are states whose density matrix satisfies
\[\hat{\rho }\ne \hat{\rho }^2\]
and a pure state otherwise.
\end{rmk}
\vspace{2ex}
\begin{proof}
Let's compare $\rho ^2$ and $\rho $
	\[{\rho }-{\rho }^2=\sum _{j}(p_{j}-p_{j}^2)|\psi _{j}\rangle \langle \psi _{j}|\]
We see that the expectation value is always bigger than 0 (sometimes, we write ${\rho }\geq {\rho }^2$). We may also ask, what is the norm of $\rho |\psi \rangle $ for an arbitrary $|\psi \rangle $? We notice that
\[\langle{\rho }\psi |{\rho }\psi \rangle =\langle \psi |{\rho }^2|\psi \rangle =\langle \psi|\rho \psi \rangle \leq  \sqrt{\langle \psi |\psi \rangle }\sqrt{\langle {\rho } \psi |{\rho } \psi\rangle } \]
Resulting in
\[\langle {\rho }\psi |{\rho }\psi \rangle\leq \langle \psi |\psi \rangle  \]
We see that the density matrix is a contraction mapping. We now prove the contraposition, that if ${\rho }^2={\rho }$, the state is pure. Last class we have seen that the eigenvalue of the matrix becomes either 1 or 0. In the spectral decomposition,
\[{\rho }=\sum _{j}p_{j}|\psi _{j}\rangle \langle \psi _{j}|\]
we see that only one of the $p_{j}$'s should become one and that the others should be zero. We have now concluded the proof. 
\end{proof}
\vspace{2ex}
\begin{thm}
(Equivalence of representations) Consider an orthonormal basis $\{|\psi _{i}\rangle \}_{i=1}^{n}$ and a not necessarily orthonormal basis $\{|\phi _{j}\rangle \}_{j=1}^{m}$. 
\[{\rho }=\sum ^{n}_{i=1}p_{i}|\psi _{i}\rangle \langle \psi _{i}|\qquad \& \qquad {\rho }=\sum ^{m}_{j=1}q_{j}|\phi   _{j}\rangle \langle \phi   _{j}|\]
As $n$ is the minimum number of vectors the span the set, $m\geq n$. Define a $m\times n$ matrix
\[V_{ij}=\sqrt{\dfrac{q_{i}}{p_{j}}}\langle \psi _{j}|\phi _{i}\rangle \]
and we see that
\[\sum ^{m}_{j=1}(V^{\dagger })_{ij}V_{jk}=\sum ^{m}_{j=1}\sqrt{\dfrac{q_{j}}{p_{i}}}\langle \phi _{j}|\psi _{i}\rangle \sqrt{\dfrac{q_{j}}{p_{k}}}\langle \psi _{k}|\phi  _{j}\rangle =\sum _{i=1}^{n}\sqrt{\dfrac{p_{i}^2}{p_{i}p_{k}}}\langle \psi _{i}|\psi _{i}\rangle \langle \psi _{k}|\psi _{i}\rangle=\delta _{ik} \]
and therefore $V^{\dagger }V=I_{n}$. After adding $m-n$ column vectors to $V$, we can make it unitary, and it would arrive at
\[\sum ^{n}_{j=1}U_{ij}\sqrt{p_{j}}|\psi _{j}\rangle =\sqrt{q_{i}}|\phi _{i}\rangle \]
This can be used to prove that $\rho $ can indeed be expressed as $\sum ^{m}_{i=1}q_{i}|\phi _{i}\rangle \langle \phi _{i}|$.
\end{thm}
\vspace{2ex}

